\documentclass[a4]{article}

% PAQUETES %%%%%%%%%%%%%%%%%%%%%%%%%%%%%%%%%%%%%%%%%%%%%%%%%%%%%%%%%%%%%%%%%%%%%
\usepackage[utf8]{inputenc}
\usepackage{framed} % Para crear cuadros alrededor del texto
\usepackage{makecell}
\usepackage[english]{babel}
%
\usepackage[fixlanguage]{babelbib}
    \bibliographystyle{IEEEtran}

\usepackage{url}

%!TEX encoding = UTF-8 Unicode
% !TEX spellcheck = es

\setlength{\textwidth}{6.5 in}
\setlength{\textheight}{8.5 in}
\setlength{\headsep}{0.5 in}

\usepackage{multirow, array} % para las tablas
\usepackage{float} % para usar [H]
\usepackage{charter}

\usepackage[breaklinks=true]{hyperref}

%
% ADD PACKAGES here:
%

\usepackage[utf8]{inputenc}
\usepackage{amsmath,longtable,siunitx, float, graphicx, dcolumn, bm, fancyhdr, authblk, subcaption, caption, hyperref, multicol, setspace}

\usepackage{bbding}
\usepackage{amssymb}
\usepackage[rightcaption]{sidecap}
%\usepackage{textcomp}

\usepackage[hmarginratio=1:1, top=32mm, columnsep=20pt, headsep=\dimexpr20mm, margin=0.7in, top=1in ]{geometry}

\usepackage{booktabs}
\usepackage{multirow}

\usepackage{algorithm}
\usepackage{enumitem}
\usepackage{tabularx}
\usepackage[noend]{algpseudocode}
\usepackage{xcolor}

\usepackage{tikz}
\usetikzlibrary{shapes.geometric, arrows.meta, positioning}

% ----------- global styles -----------
\tikzset{
  font=\scriptsize,                      % tiny text everywhere
  every node/.style       = {transform shape},   % respect scale
  startstop/.style        = {ellipse, draw, fill=blue!10,
                             inner sep=1pt, minimum width=13mm},
  process/.style          = {rectangle, draw, fill=orange!15,
                             text width=26mm, inner sep=2pt, align=center},
  decision/.style         = {diamond, draw, fill=green!15, aspect=2.3,
                             text width=20mm, inner sep=1pt, align=center},
  arrow/.style            = {->, >=Stealth, shorten >=1pt, shorten <=1pt}
}

\newcommand{\TODO}{\textbf{\textcolor{red}{TODO}}}

\renewcommand{\algorithmicrequire}{\textbf{Input:}}
\renewcommand{\algorithmicensure}{\textbf{Output:}}

\lhead{\includegraphics[scale=0.4]{CETC.png}}
\rhead{\includegraphics[scale=0.025]{Cornell-University-Logo.png}}

\newtheorem{theo}{Theorem}
\newenvironment{theorem}
  {\begin{mdframed}\begin{theo}}
  {\end{theo}\end{mdframed}}
  
\newtheorem{prop}{Proposición}
\newenvironment{proposition}
  {\begin{mdframed}\begin{prop}}
  {\end{prop}\end{mdframed}}

\newtheorem{lem}{Lema}
\newenvironment{lema}
  {\begin{mdframed}\begin{lem}}
  {\end{lem}\end{mdframed}}
  
\newtheorem{defi}{Definition}
\newenvironment{definition}
  {\begin{mdframed}\begin{def}}
  {\end{def}\end{mdframed}}
  
\newtheorem{proof}{Proof}
\newtheorem{remark}{Remark}

\usepackage[numbers]{natbib}
\renewcommand{\thesection}{\Roman{section}}
%\setlength{\parskip}{\baselineskip}%

\begin{document}

\thispagestyle{fancy}

\begin{center}
    \textbf{\large Orbit Odyssey - Computer Vision} \\
    {Field Definition}
    
    \vspace{0.7em}

    David Alejandro Fuquen \\

    {\small \today}

\end{center}

\section*{Playing Field Definition}

\subsection*{Summary}
This part of Phase 1 defines all the structure, positions, conditions and specifications for the playing field. This is the first part of the competition manual. For this case, we will consider the partition of the original field will be taken away. 


\section*{Field dimensions, boundaries and markings}
We will use the same dimensions and boundaries as the original Space Odyssey. The same 4x8 field will be used, without the inner field partition. 

The playing field is a rectangular arena enclosed by low-profile PVC pipe barriers on all four sides.

\begin{table}[H]
\centering
\begin{tabular}{ll}
\toprule
\textbf{Property} & \textbf{Value} \\
\midrule
Total length (X-axis) & 254.00 cm (100 in) \\
Total width (Y-axis) & 142.24 cm (56 in) \\
Total floor area & 3.61 m\textsuperscript{2} \\
Perimeter walls & Low-profile PVC pipe barriers \\
Corner zones & 30.48 $\times$ 30.48 cm (12 $\times$ 12 in) at each corner \\
Partition & \textbf{Removed} \\
\bottomrule
\end{tabular}
\caption{Field dimensions and boundaries.}
\end{table}

The coordinate system places the origin $(0, 0)$ at the bottom-left corner of the field. The X-axis runs left to right (0 to 254 cm) and the Y-axis runs bottom to top (0 to 142.24 cm).

//TODO (David): Add an image with the field. DON'T DO THIS, I'll do it myself. 

\section*{Robot starting zone and placement rules}
The robot will start in the center of the field. It will be surrounded by the 4 objects, and the team will decide the initial orientation. This will be the only human-made decision of the challenge, the only part in which the team has any kind of decision-making abilities.
\section*{Object quantity, placement templates, and clearance constraints}
\subsection*{Object count and class breakdown}
This layout uses \textbf{3 objects}: 1 Rocket (contact) + 1 Tire (avoid) + 1 Tin Can (avoid). Each object has an approximate radius of 6 cm. This is a \textbf{center-start fallback} layout designed for the case where \texttt{main.py} cannot be modified --- the robot spins in place at the center of the field and detects surrounding objects without significant translational movement.

\subsection*{Robot starting position}

The robot is placed at the center of the field at \textbf{(50, 28) in = (127.00, 71.12) cm}. The team selects the initial heading.

\subsection*{Object positions}

The three objects are arranged in an approximately equilateral triangle around the robot, at $\sim$50 cm from the center.

\begin{table}[H]
\centering
\begin{tabular}{clcccc}
\toprule
\textbf{ID} & \textbf{Location} & \textbf{X (in)} & \textbf{Y (in)} & \textbf{X (cm)} & \textbf{Y (cm)} \\
\midrule
P1 & Top-center & 50 & 48 & 127.00 & 121.92 \\
P2 & Lower-left & 33 & 18 & 83.82 & 45.72 \\
P3 & Lower-right & 67 & 18 & 170.18 & 45.72 \\
\bottomrule
\end{tabular}
\caption{Object positions --- Layout NP\_4 (center start).}
\end{table}

\subsection*{Distance from robot to each object}

\begin{table}[H]
\centering
\begin{tabular}{ccc}
\toprule
\textbf{ID} & \textbf{Distance to robot (cm)} & \textbf{Distance to robot (in)} \\
\midrule
P1 & 50.8 & 20.0 \\
P2 & 50.0 & 19.7 \\
P3 & 50.0 & 19.7 \\
\bottomrule
\end{tabular}
\caption{Distance from robot center to each object.}
\end{table}

\subsection*{Clearance verification}

\textbf{Minimum clearances:} object center to wall $\geq$ 20 cm, object center to object center $\geq$ 35 cm, objects must be outside all corner zones.

\begin{table}[H]
\centering
\begin{tabular}{ccccccc}
\toprule
\textbf{ID} & \textbf{Left wall} & \textbf{Right wall} & \textbf{Bottom wall} & \textbf{Top wall} & \textbf{Min wall} & \textbf{Nearest obj.} \\
\midrule
P1 & 127.0 & 127.0 & 121.9 & 20.3 & 20.3 \checkmark & 88.0 (P2) \checkmark \\
P2 & 83.8 & 170.2 & 45.7 & 96.5 & 45.7 \checkmark & 86.4 (P3) \checkmark \\
P3 & 170.2 & 83.8 & 45.7 & 96.5 & 45.7 \checkmark & 86.4 (P2) \checkmark \\
\bottomrule
\end{tabular}
\caption{Clearance verification (all values in cm).}
\end{table}

\subsection*{Suggested class rotation templates}

With only 3 objects and 3 classes, there are 6 possible permutations. Three representative templates:

\begin{table}[H]
\centering
\begin{tabular}{lccc}
\toprule
\textbf{Template} & \textbf{P1} & \textbf{P2} & \textbf{P3} \\
\midrule
A & Rocket & Tire & Tin Can \\
B & Tin Can & Rocket & Tire \\
C & Tire & Tin Can & Rocket \\
\bottomrule
\end{tabular}
\caption{Class rotation templates.}
\end{table}






\end{document}
